\section{Introduction}
Large language models (LLMs) have recently demonstrated strong multi-step reasoning capabilities when prompted to generate explicit chains of thought (CoT), in which intermediate reasoning steps are expressed as natural language tokens prior to producing a final answer \cite{wei2022chain}. While effective, this paradigm tightly couples reasoning to the discrete token space of language, forcing models to externalize internal computation through sequential text generation. Language tokens, however, are optimized for communication rather than for the core requirements of reasoning such as abstraction, conceptualization, and computation \cite{fedorenko2024language,amalric2019distinct,monti2007functional,monti2009boundaries,monti2012thought,fedorenko2011functional}. Moreover, autoregressive decoding commits to a single partial trajectory token by token, and the computational cost of reasoning scales linearly with the number of generated tokens \cite{saparov2022language}. These constraints make detailed reasoning both expensive and inherently sequential.

A substantial body of work has sought to alleviate these issues through improved prompting strategies and decoding-time search. Techniques such as decomposed prompting \cite{khot2022decomposed}, least-to-most prompting \cite{zhou2022least}, self-consistency and guided beam-style decoding \cite{xie2023self}, and tree-based search over reasoning trajectories \cite{yao2023tree} improve robustness by exploring or aggregating multiple textual reasoning paths. Related approaches introduce test-time deliberation signals that encourage models to ``think'' before committing to an answer \cite{zelikman2024quiet}. Despite their success, these methods operate entirely in token space and therefore inherit its fundamental limitation: each reasoning step must still be generated, evaluated, and committed sequentially. As a result, they improve exploration within the same substrate rather than enabling more abstract or information-dense internal computation.

Recent work has therefore explored shifting reasoning from explicit language tokens into the model’s continuous hidden state space. \emph{Chain of Continuous Thought} (Coconut) introduced the idea of replacing segments of the textual chain-of-thought with continuous latent representations that are never decoded as language and are trained implicitly through their effect on final answer likelihood \cite{hao2024training}. Empirically, Coconut demonstrates that models trained in this manner can exhibit emergent search behaviors resembling breadth-first exploration while requiring fewer explicit reasoning tokens. Complementing this empirical evidence, theoretical analyses show that continuous latent representations can be strictly more expressive than discrete token sequences: \emph{Reasoning by Superposition} proves that a single latent step can encode superpositions of multiple candidate reasoning states that would require sequential exploration in token space \cite{zhu2025reasoning}. Together, these results establish continuous latent reasoning as a promising alternative substrate for complex reasoning.

Despite this promise, existing methods of continuous latent reasoning suffer from a critical structural limitation. In Coconut, continuous thoughts are generated \emph{sequentially}: each latent vector is produced by feeding the previous hidden state back as the next input embedding \cite{hao2024training}. This design introduces a recurrent unrolling within the Transformer, reinstating an RNN-like dependency structure and causing the critical path length to grow linearly with the number of latent steps. Consequently, training becomes substantially slower, hardware parallelism is poorly utilized, and optimization stability degrades as latent depth or model size increases, as gradients must propagate through a long sequence of dependent latent states, causing compounding of early representational instabilities.

Importantly, this recurrent structure imposes not only a computational bottleneck but also a representational one. When latent thoughts are generated by recursively transforming prior latent states, each latent must accumulate and carry forward information from earlier steps, as future computation has access only to the most recent latent. The absence of an explicit stage boundary encourages successive latents to entangle multiple reasoning roles within a single evolving representation, rather than forming distinct, stage-specific abstractions. Moreover, because learning signals are applied only through the final answer likelihood, recurrent unrolling biases the model toward encoding global summaries early rather than organizing reasoning into separable stages. In contrast, reasoning architectures that re-synthesize latent representations from the current context at each stage can refactor information anew, discouraging unstructured accumulation and enabling specialization across stages. We examine these representational effects empirically through detailed analyses later in the paper.

Several recent methods aim to improve the efficiency of latent reasoning but do not address this architectural bottleneck. CoLaR dynamically compresses chains of thought into fewer latent steps and optimizes reasoning length via reinforcement learning, but it retains autoregressive latent generation and explicitly supervises latents to preserve token-level semantics \cite{tan2025think}. CODI similarly distills explicit CoT into continuous latent representations through self-distillation, treating latents as compact surrogates for textual reasoning traces \cite{shen2025codi}. Token Assorted mixes latent and text tokens during reasoning but likewise preserves sequential latent computation \cite{su2025token}. These approaches reduce the number of latent steps or compress existing reasoning traces, yet they do not remove the core sequential dependency induced by recurrent latent unrolling, nor do they enable latent representations to reorganize computation across reasoning stages.

In this work, we introduce \emph{Parallel Continuous Thought} (PaCT), a reformulation of continuous latent reasoning that eliminates recurrent latent unrolling altogether.
\textit{PaCT challenges the implicit assumption that continuous latent reasoning must rely on latent state propagation across steps, showing instead that reasoning depth can be achieved through repeated, context-grounded latent resynthesis by applying parallelism.}
Instead of generating latent thoughts sequentially, PaCT synthesizes a set of latent representations \emph{in parallel} at each reasoning stage using a learnable codebook of query vectors and cross-attention to the current textual context. These latents are then refined through bidirectional interaction within a latent representation set before conditioning subsequent decoding. By re-synthesizing latents at each stage rather than recursively accumulating prior latent states, PaCT decouples reasoning capacity from serial computation depth, enables efficient hardware parallelism, and significantly reduces training and optimization overhead, while encouraging stage-wise specialization of latent representations and preserving standard causal decoding.

Empirically, PaCT matches or exceeds the reasoning accuracy of Coconut across diverse mathematical and logical reasoning benchmarks and model scales, while achieving $4$--$9\times$ reductions in training time (Figure~\ref{fig:concept_diagram_intro}). Beyond efficiency, our analyses show that PaCT produces information-rich latent reasoning where fewer latent states suffice to achieve strong performance, increasing reasoning depth yields consistent gains without increasing latent width, and successive stages encode distinct, context-grounded information rather than redundantly accumulating latent history. Together, these results demonstrate that parallel latent synthesis is not merely an acceleration strategy, but an architectural improvement that enables both scalable and higher-quality continuous latent reasoning. Our main contributions are -\\
\textbf{(1) Stage-wise latent resynthesis without recurrent unrolling.}
We eliminate recurrent latent state propagation by re-synthesizing latent reasoning states at each stage from the current context, decoupling reasoning depth from serial computation. This enables parallel latent synthesis via cross-attention and yields substantially faster training while matching or improving reasoning accuracy (Figure~\ref{fig:pact_overview}).\\\\
\textbf{(2) Information-rich latents via within-stage deliberation.}
We introduce bidirectional interaction among latent states within each reasoning stage, producing more information-dense latent representations and improving reasoning capacity without increasing latent width (Figure~\ref{fig:latent_attention}).
% Following are the main contributions of our work-
% \vspace{-6mm}
% \begin{itemize} \itemsep0pt \parsep0pt \topsep2pt
%     \item \textbf{Stage-wise latent resynthesis without recurrent unrolling.} We eliminate recurrent latent state propagation by re-synthesizing latent reasoning states at each stage from the current context, decoupling reasoning depth from serial computation. This enables parallel latent synthesis via cross-attention and yields substantially faster training while matching or improving reasoning accuracy (Figure~\ref{fig:pact_overview}).
%     \item \textbf{Information-rich latents via within-stage deliberation.} We introduce bidirectional interaction among latent states within each reasoning stage, producing more information-dense latent representations and improving reasoning capacity without increasing latent width (Figure~\ref{fig:latent_attention}).
% \end{itemize}
